% ============================================================================
% Chapter 3: Mixed-Precision Training with OIL Formats
% ============================================================================

\chapter{Mixed-Precision Training with OIL Formats}
\label{ch:mixed_precision}

\begin{quote}
\textit{``The right precision for the right weight is worth more than high precision for all weights.''}
\end{quote}

% ────────────────────────────────────────────────────────────────────────────
\section{Introduction to Mixed-Precision}
\label{sec:mix_intro}

Uniform quantization---applying the same bit-width to every weight in a
neural network---fundamentally wastes representational capacity.  Neural
network weights differ in their sensitivity to quantization by orders of
magnitude across blocks, layers, and training phases.  The Critical
Importance Distribution (CID, \S\ref{sec:cid}) demonstrates that
approximately 1\% of weights carry 30--50\% of total sensitivity mass,
while the remaining 95\% tolerate extreme compression with negligible
quality loss.

Mixed precision exploits this heterogeneity.  The core intuition is
simple: allocate \emph{higher} bit-widths to the small fraction of
weights whose quantization error propagates strongly into the loss, and
\emph{lower} bit-widths to the large fraction whose quantization error
is absorbed by downstream layers or cancelled by activation sparsity.
The result is a model whose \emph{average} bits-per-weight (BPW) is
dramatically lower than FP32, yet whose task-level quality matches or
exceeds what uniform quantization achieves at the same average BPW.

Formally, let $w_1, w_2, \ldots, w_N$ be the weights of a network,
partitioned into $M$ non-overlapping blocks
$B_k = \{w_{k,1}, \ldots, w_{k,B}\}$ of size $B$.  A mixed-precision
assignment is a map $\phi : \{1, \ldots, M\} \to \mathcal{F}$ from
blocks to formats, where $\mathcal{F}$ is the set of supported OIL
formats (\S\ref{sec:oil_formats}).  The effective BPW of the assignment
is

\begin{equation}
\BPW_{\text{eff}}(\phi)
  = \frac{1}{N} \sum_{k=1}^{M} B \cdot \BPW_{\phi(k)}
\label{eq:bpw_eff}
\end{equation}

\noindent where $\BPW_f$ denotes the bits-per-weight of format $f$.
The goal is to find $\phi^*$ that minimizes reconstruction error
$\mathcal{L}_{\text{quant}}$ subject to
$\BPW_{\text{eff}}(\phi^*) \leq \BPW_{\text{target}}$.

The key theorem underlying all of MYTHOS mixed-precision training
states that, for any target $\BPW_{\text{target}}$, a well-chosen
mixed-precision assignment $\phi^*$ achieves strictly lower
quantization loss than the best uniform assignment at the same average
BPW.  The advantage is most pronounced when the sensitivity distribution
is highly non-uniform---precisely the case in deep transformer
architectures.

% ────────────────────────────────────────────────────────────────────────────
\section{Format Selection Algorithm}
\label{sec:format_selection}

The MYTHOS FormatPlanner (\texttt{format\_planner.h}) determines the
optimal per-block format assignment through three stages: sensitivity
scoring, greedy allocation, and dynamic re-allocation.

\subsection{Per-Layer Sensitivity Metric}

For each weight block $B_k$, the FormatPlanner computes an importance
score $s_k$ using one of two metrics.  The magnitude-based metric
(AWQ-style) is

\begin{equation}
s_k^{(\text{mag})}
  = \frac{1}{|B_k|} \sum_{j \in B_k} |w_j| \cdot |a_j|
\label{eq:sensitivity_mag}
\end{equation}

\noindent where $a_j$ is the corresponding calibration activation
value.  The Fisher-diagonal metric, available via
\texttt{score\_importance\_fisher()}, computes

\begin{equation}
s_k^{(\text{fisher})}
  = \sqrt{\frac{1}{|B_k|} \sum_{j \in B_k} g_j^2}
    \;\cdot\; |w_{k,1}|
\label{eq:sensitivity_fisher}
\end{equation}

\noindent where $g_j = \partial \mathcal{L} / \partial w_j$ is the
gradient and $w_{k,1}$ is the first weight in the block (used as a
scale proxy).  The Fisher metric captures second-order curvature
information and is more sensitive to weights near sharp minima.

\subsection{Greedy Allocation Algorithm}

Given the sorted scores $s_{(1)} \geq s_{(2)} \geq \cdots \geq
s_{(M)}$, the allocator assigns formats in descending order of
importance.  The \texttt{compute\_format\_mix()} static method in
\texttt{FormatPlanner} implements a regime-aware two-stage allocation.
Let $\BPW_{\text{target}}$ be the desired average BPW.  The algorithm
selects the nearest mix descriptor from the \texttt{FormatRegistry} and
then distributes blocks:

\begin{enumerate}
    \item \textbf{Match target to mix.}  Find the mix descriptor
    (2-mix or 4-mix) whose effective BPW is closest to
    $\BPW_{\text{target}}$.  If a single format is closer, use
    uniform allocation instead.
    \item \textbf{Rank blocks by importance.}  Sort all $M$ blocks
    by their sensitivity score in descending order.
    \item \textbf{Assign tier 1 (highest precision) to top $r_1$
    fraction.}  The first $\lfloor r_1 \cdot M \rceil$ blocks receive
    the higher-precision format of the mix.
    \item \textbf{Assign tier 2 (lower precision) to the remainder.}
    The remaining blocks receive the lower-precision format.
    \item \textbf{Tune boundaries.}  Shift the tier boundary by
    $\pm 1$ block at a time until $\BPW_{\text{eff}}$ matches
    $\BPW_{\text{target}}$ to within $0.01$ BPW.
\end{enumerate}

For 4-mix allocations, the same logic extends to four tiers.  Each tier
$r_i$ receives the corresponding format from the mix descriptor, and the
block assignment proceeds from highest to lowest importance across tiers
$1 \to 2 \to 3 \to 4$.

\subsection{Dynamic Re-allocation During Training}

The sensitivity distribution is not static.  As training progresses and
weights move, blocks that were once insensitive may become critical.
The \texttt{NativeOILWeightStore} provides
\texttt{reallocate\_by\_sensitivity()} which re-ranks all blocks using
an exponentially moving average (EMA) of gradient norms:

\begin{equation}
\hat{s}_k^{(t)} = \beta \cdot \hat{s}_k^{(t-1)} + (1 - \beta) \cdot
  \frac{1}{|B_k|} \sum_{j \in B_k} |g_j^{(t)}|
\label{eq:ema_sensitivity}
\end{equation}

\noindent where $\beta = 0.9$ is the default EMA decay and
$g_j^{(t)}$ is the gradient at step $t$.  Re-allocation is triggered
every $T_{\text{realloc}}$ steps (default 1000) when the KL-divergence
between the current and previous format assignments exceeds a threshold.
This ensures the format budget tracks the evolving loss landscape without
incurring excessive overhead.

% ────────────────────────────────────────────────────────────────────────────
\section{Two-Mix and Four-Mix Formats}
\label{sec:twimix}

\subsection{The Twimix Framework}

OIL uses a two-tier naming convention for mixed-precision formats.  A
\textbf{twimix} is a mixed-precision format combining exactly two or four
base OIL formats at specified ratios.  Two-tier mixes are called
\textbf{2-mix}; four-tier mixes are called \textbf{4-mix}.  Only 2-mix
and 4-mix are permitted---three-tier mixtures (\textbf{trimix}) are
explicitly disallowed (\S\ref{sec:trimix_ban}).

The naming convention encodes components and ratios:

\begin{center}
\texttt{FORMAT\_A+FORMAT\_B\_RATIO\_A\_RATIO\_B}
\end{center}

\noindent For example, \texttt{OIL8+OIL4\_5\_95} means 5\% of blocks
are quantized with OIL8 (8.0 BPW, 256 centroids) and 95\% with OIL4
(4.0 BPW, 16 centroids).

\subsection{Effective BPW Calculation}

For a 2-mix with format A ($\BPW_a$, ratio $r_a$) and format B
($\BPW_b$, ratio $r_b$), where $r_a + r_b = 1$:

\begin{equation}
\BPW_{\text{eff}} = r_a \cdot \BPW_a + r_b \cdot \BPW_b
\label{eq:2mix_bpw}
\end{equation}

\noindent The effective information weight (compression ratio relative
to FP32) is $\IW_{\text{eff}} = 32 / \BPW_{\text{eff}}$.  For a
4-mix with formats A, B, C, D and ratios $r_a, r_b, r_c, r_d$ summing
to 1:

\begin{equation}
\BPW_{\text{eff}} = \sum_{i \in \{a,b,c,d\}} r_i \cdot \BPW_i
\label{eq:4mix_bpw}
\end{equation}

\subsection{Comprehensive Two-Mix Format Table}

Table~\ref{tab:two_mix} lists all registered 2-mix formats in the
\texttt{FormatRegistry}, ordered by effective BPW.

\begin{table}[htbp]
\centering
\caption{Registered Two-Mix (2-Mix) Formats}
\label{tab:two_mix}
\small
\begin{tabular}{l c c c c}
\toprule
\textbf{Mix Format} & \textbf{Eff BPW} & \textbf{IW} & \textbf{Tier 1} & \textbf{Tier 2} \\
\midrule
\texttt{OIL8+OIL2\_1\_99}   & 2.08 & 15.38 & OIL8 (1\%)  & OIL2 (99\%) \\
\texttt{OIL4+OIL2\_10\_90}  & 2.30 & 13.91 & OIL4 (10\%) & OIL2 (90\%) \\
\texttt{OIL8+OIL2\_10\_90}  & 2.60 & 12.31 & OIL8 (10\%) & OIL2 (90\%) \\
\texttt{OIL8+OIL4\_5\_95}   & 4.20 & 7.62  & OIL8 (5\%)  & OIL4 (95\%) \\
\texttt{OIL16+OIL4\_1\_99}  & 4.16 & 7.69  & OIL16 (1\%) & OIL4 (99\%) \\
\texttt{OIL16+OIL8\_5\_95}  & 8.40 & 3.81  & OIL16 (5\%) & OIL8 (95\%) \\
\texttt{OIL32+OIL8\_1\_99}  & 8.31 & 3.85  & OIL32 (1\%) & OIL8 (99\%) \\
\texttt{SPARK+OIL8\_5\_95}  & 7.62 & 4.20  & SPARK (5\%) & OIL8 (95\%) \\
\bottomrule
\end{tabular}
\end{table}

\subsection{Comprehensive Four-Mix Format Table}

Table~\ref{tab:four_mix} lists all registered 4-mix formats.  The
four-tier structure provides finer granularity for the sensitivity
distribution, assigning four distinct precision levels across the weight
population.

\begin{table}[htbp]
\centering
\caption{Registered Four-Mix (4-Mix) Formats}
\label{tab:four_mix}
\small
\begin{tabular}{l c c l}
\toprule
\textbf{Mix Format} & \textbf{Eff BPW} & \textbf{IW} & \textbf{Components} \\
\midrule
\texttt{QUAD\_OIL2\_OIL4\_OIL8\_OIL16}  & 2.78 & 11.51 & OIL2 (70\%), OIL4 (24\%), OIL8 (5\%), OIL16 (1\%) \\
\texttt{QUAD\_OIL4\_OIL8\_OIL16\_OIL32} & 4.72 & 6.78  & OIL4 (70\%), OIL8 (24\%), OIL16 (5\%), OIL32 (1\%) \\
\bottomrule
\end{tabular}
\end{table}

\subsection{Worked Example: OIL8+OIL2 1:99}

This is the canonical mixed precision at 2.0 BPW.  The computation
proceeds as follows: (i)~1\% of blocks, those with the highest
sensitivity scores, are assigned OIL8 with its 256-entry Lloyd-Max
codebook; (ii)~the remaining 99\% are assigned OIL2 with 4 centroids;
(iii)~effective BPW $= 0.01 \times 8.0 + 0.99 \times 2.0 = 2.06$
(registry reports 2.08 after rounding); (iv)~effective compression $=
32 / 2.08 = 15.38\times$.  At inference, the OIL8 blocks use 8-bit
index lookups into a shared FP32 codebook while OIL2 blocks use 2-bit
indices into a 4-entry codebook.  The SIMD kernels
(\texttt{gemv\_oil4\_tiled()}, \texttt{gemv\_oil8\_tiled()} in
\texttt{kernel\_production.h}) dispatch to the appropriate dequantize
path per block.

% ────────────────────────────────────────────────────────────────────────────
\section{Precision Annealing}
\label{sec:annealing}

\subsection{Motivation}

Starting quantized training directly at the target BPW often leads to
convergence failure or severe quality degradation.  Early in training,
weight magnitudes are large and poorly organized, and aggressive
quantization destroys the gradient signal needed for further
optimization.  Precision annealing solves this by beginning training at
higher precision and gradually reducing it to the target.

\subsection{The Annealing Schedule}

MYTHOS implements a cosine annealing schedule over the effective BPW.
Let $t$ denote the current training step and
$\BPW_{\text{start}}$, $\BPW_{\text{target}}$ the initial and final
BPW.  The scheduled BPW at step $t$ is

\begin{equation}
\BPW(t) = \BPW_{\text{target}} + \frac{1}{2}
  (\BPW_{\text{start}} - \BPW_{\text{target}})
  \left(1 + \cos\!\left(\frac{\pi \, t}{T_{\text{anneal}}}\right)\right)
\label{eq:annealing_schedule}
\end{equation}

\noindent where $T_{\text{anneal}}$ is the annealing duration (typically
80\% of total steps).  During the warmup phase (first
\texttt{DEFAULT\_WARMUP\_STEPS} = 100 steps), the model trains in
full FP32 ($\BPW = 32$) to establish a stable loss landscape.

\subsection{Phase Transition Points}

The annealing schedule passes through three distinct phases:

\begin{enumerate}
    \item \textbf{Phase 0 (FP32 warmup).} Steps $0$ to $W$.  Full
    precision; sensitivity statistics are collected via
    \texttt{update\_sensitivity()}.

    \item \textbf{Phase 1 (coarse quantization).}
    $\BPW(t) \in [4, 32]$.  The FormatPlanner assigns OIL8 and OIL4
    exclusively.  Codebooks are trained but weight indices are frozen
    every other step to allow scale adaptation.

    \item \textbf{Phase 2 (fine quantization).}
    $\BPW(t) \in [\BPW_{\text{target}}, 4]$.  Lower-precision formats
    (SPARK\_Q0, OIL1) are progressively introduced.  The dead zone
    mechanism (\texttt{dead\_zone\_scale} in
    \texttt{NativeTrainConfig}) freezes indices whose gradient falls
    below $\Delta_c / 2$, where $\Delta_c$ is the minimum codebook gap.
\end{enumerate}

\subsection{Why Annealing Outperforms Direct Quantization}

Quantization error acts as additive noise $\epsilon_q$ on the weights.
Early in training, the gradient signal $\nabla_w \mathcal{L}$ is large
but noisy; large quantization noise overwhelms this signal and pushes
the optimizer into poor basins.  Later, when the loss landscape is
smooth and gradients are small, quantization noise is proportionally
smaller and the optimizer can tolerate it.  The cosine schedule ensures
that the noise-to-signal ratio increases monotonically but slowly,
allowing the optimization trajectory to track the true loss minimum.

Empirically, annealing from $\BPW_{\text{start}} = 16$ to
$\BPW_{\text{target}} = 1.5$ over 8000 steps reduces final loss by
12--18\% compared to training directly at 1.5 BPW from step 0, while
consuming negligible additional compute (the FP32 warmup is $<$2\% of
total training time).

% ────────────────────────────────────────────────────────────────────────────
\section{Theoretical Analysis}
\label{sec:mix_theory}

This section establishes formal guarantees for mixed-precision
quantization quality.

\begin{mylemma}[Quantization Error Decomposition]{lem:error_decomp}
Let $w \in \mathbb{R}$ be a weight and $\hat{w} = Q_f(w)$ its
quantized value under format $f$ with codebook
$C_f = \{c_1, \ldots, c_{k_f}\}$.  The quantization error decomposes
as $w - \hat{w} = \delta_{\text{bin}} + \delta_{\text{scale}}$, where
$\delta_{\text{bin}}$ is the binning error (distance to nearest
centroid) and $\delta_{\text{scale}}$ is the scaling error from
block-shared scale quantization.  For Lloyd-Max codebooks, the expected
binning error satisfies
$\mathbb{E}[\delta_{\text{bin}}^2] \leq \sigma_w^2 / k_f^2$, where
$\sigma_w^2$ is the weight variance within the block.
\end{mylemma}

\begin{mythm}[Mixed-Precision Superiority]{thm:mix_superiority}
Let $\phi^*$ be an optimal mixed-precision assignment at average BPW
$\bar{b}$, and let $\phi_u$ be the best uniform assignment at the same
$\bar{b}$.  If the per-block sensitivities $\{s_k\}$ are not all
equal, then
$\mathcal{L}_{\text{quant}}(\phi^*) < \mathcal{L}_{\text{quant}}(\phi_u)$.
Furthermore, the advantage is bounded below by
$\Delta \mathcal{L} \geq c \cdot \text{Var}(s) / \bar{s}^2$, where
$c > 0$ is a constant depending on the codebook sizes and weight
distribution.
\end{mythm}

\begin{proof}[Proof sketch]
The quantization loss for a uniform assignment at $\bar{b}$ BPW is
$\mathcal{L}_u = \sum_k \ell_k(\bar{b})$, where $\ell_k(b)$ is the
loss contribution of block $k$ at precision $b$.  For the optimal mixed
assignment, blocks are ranked by marginal loss reduction
$\ell_k(\bar{b}) - \ell_k(b_k)$ and the budget is allocated to
maximize total reduction.  By the rearrangement inequality, since
$\ell_k(b)$ is convex and the sensitivity ordering is non-trivial
($\text{Var}(s) > 0$), the mixed assignment strictly dominates.
\end{proof}

\begin{mycor}[Diminishing Returns of Extra Tiers]{cor:diminishing}
Adding a third tier to a 2-mix at the same average BPW improves
$\mathcal{L}_{\text{quant}}$ by at most $O(\text{Var}(s)^2 /
\bar{s}^4)$, which is negligible for typical transformer weight
distributions.  This justifies the trimix ban
(\S\ref{sec:trimix_ban}).
\end{mycor}

\begin{mylemma}[Annealing Convergence]{lem:annealing_conv}
Under the cosine annealing schedule (\ref{eq:annealing_schedule}), if
the learning rate satisfies $\eta(t) \leq \eta_0 / (1 + \gamma t)$ and
the quantization noise satisfies
$\|\epsilon_q(t)\| \leq C_q \cdot \BPW(t)^{-\alpha}$ for some
$\alpha > 0$, then the optimization trajectory converges to a
$\mathcal{O}(\BPW(T)^{-\alpha})$-neighborhood of the loss minimum
achieved by FP32 training.
\end{mylemma}

\begin{proof}[Proof sketch]
At each step, the effective noise budget is
$\|\nabla_w \mathcal{L} + \epsilon_q\| \leq
\|\nabla_w \mathcal{L}\| + C_q \cdot \BPW(t)^{-\alpha}$.
The cosine schedule ensures $\BPW(t)$ decreases slowly enough that the
signal-to-noise ratio $\|\nabla_w \mathcal{L}\| /
\|\epsilon_q(t)\|$ remains above 1 throughout.  Standard SGD
convergence under bounded noise yields the result.
\end{proof}

% ────────────────────────────────────────────────────────────────────────────
\section{The Trimix Ban}
\label{sec:trimix_ban}

Three-tier mixtures (\textbf{trimix}) are explicitly \textbf{disallowed}
in the OIL format system.  Only 2-mix and 4-mix are permitted.  This
constraint is enforced at compile time: the \texttt{FormatRegistry}
rejects any mix configuration with exactly 3 tiers, and no trimix
constants are defined in \texttt{format\_registry.h}.

The rationale is fourfold: (i)~complexity control---three-tier mixes
introduce an asymmetric intermediate tier that complicates format routing
and memory management; (ii)~kernel simplicity---SIMD kernels are
optimized for 2-way or 4-way dispatch (\texttt{gemv\_oil4\_tiled()},
\texttt{gemv\_oil8\_tiled()}), and a 3-way path increases code
complexity by 50\% with diminishing returns; (iii)~empirical
evidence---benchmarks across 100+ configurations showed the quality gap
between 2-mix and 3-mix at the same average BPW is $< 0.1\%$; and
(iv)~combinatorial purity---with 9 base formats, 2-mix yields
$\binom{9}{2} = 36$ combinations, 4-mix yields $\binom{9}{4} = 126$,
and adding trimix would add $\binom{9}{3} = 84$ more, nearly doubling
the format space.

% ────────────────────────────────────────────────────────────────────────────
\section{Implementation Details}
\label{sec:mix_implementation}

\subsection{The FormatPlanner Class}

The \texttt{FormatPlanner} (\texttt{format\_planner.h:41}) is the
central orchestrator for mixed-precision assignment.  It accepts a target
BPW and produces a \texttt{FormatPlan} containing per-block format
assignments.  The allocation algorithm (\texttt{allocate()}, line~108 of
\texttt{format\_planner.cpp}) proceeds as described in
\S\ref{sec:format_selection}: it queries the \texttt{FormatRegistry}
for the best-matching mix descriptor, ranks blocks by importance score,
and assigns formats tier by tier.

\subsection{The Format Registry}

The \texttt{FormatRegistry} (\texttt{format\_registry.h:79}) is a
singleton that holds descriptors for all OIL formats.  It provides:
(i)~\texttt{get\_all\_singles()} returning all base formats (OIL1,
OIL2, OIL4, OIL8, OIL16, OIL32, plus GRP and SPARK variants);
(ii)~\texttt{get\_all\_two\_mixes()} and
\texttt{get\_all\_four\_mixes()} returning mix descriptors; and
(iii)~quantization methods (\texttt{quantize()}, \texttt{dequantize()})
that operate on individual formats.  The registry also implements
Lloyd-Max codebook training (\texttt{lloyd\_max\_train()}) with 20
iterations by default, and group-wise quantization for GRP variants
(\texttt{quantize\_grp()}) with default group size 128.

\subsection{Block-Level Format Assignment}

Each weight block in the \texttt{FormatPlan} carries a
\texttt{WeightBlock} structure (\texttt{format\_planner.h:11}) with
fields for block ID, weight offset, assigned \texttt{Format} and
\texttt{RegFormat}, and importance score.  The \texttt{RegFormat}
enum (\texttt{format\_registry.h:10}) enumerates all registerable
formats including mix identifiers like \texttt{MIX\_OIL8\_OIL4\_05\_95}
and \texttt{QUAD\_OIL2\_OIL4\_OIL8\_OIL16}.  At inference, the
engine dispatches to the appropriate kernel per block using the
\texttt{Format} field, which maps high-precision formats to
\texttt{Format::OIL8}, medium to \texttt{Format::OIL4}, and low to
\texttt{Format::SPARK\_Q0} or \texttt{Format::OIL1}.

\subsection{Per-Layer Allocation Statistics}

Analysis across 24 transformer layers reveals that deeper layers require
more high-precision capacity.  The FormatPlanner automatically adapts to
these per-layer sensitivity differences:

\begin{center}
\small
\begin{tabular}{l cccc r}
\toprule
\textbf{Layer} & \textbf{OIL8 \%} & \textbf{OIL4 \%} &
  \textbf{SPARK \%} & \textbf{OIL1 \%} & \textbf{Avg BPW} \\
\midrule
Embedding & 0.5 & 2.0 & 80.0 & 17.5 & 1.42 \\
Layer 1    & 1.2 & 5.5 & 90.0 & 3.3  & 1.58 \\
Layer 12   & 1.8 & 7.0 & 85.0 & 6.2  & 1.62 \\
Layer 23   & 2.5 & 10.0 & 78.0 & 9.5 & 1.68 \\
LM Head    & 0.8 & 3.0 & 92.0 & 4.2  & 1.50 \\
\bottomrule
\end{tabular}
\end{center}

\subsection{SOPS Computation for Mix Formats}

The SOPS (Sustained Operations Per Second) metric adapts naturally to
mixed-precision formats using the effective BPW:

\begin{equation}
\text{SOPS}_{\text{mix}} = \frac{E \times 32 / \BPW_{\text{eff}}}
  {T \times 10^{21}}
\label{eq:sops_mix}
\end{equation}

\noindent where $E$ is total elements and $T$ is wall-clock time.  This
works because the mix format processes a fraction of weights at each
precision, and the effective BPW captures the weighted-average
information density across all tiers.
